\chapter{Sets, Relations, and Functions}

This chapter covers the foundational building blocks of mathematical language: number systems, sets, relations, and functions. While these topics might seem entirely abstract at first glance, they are the very same tools used to formalize data structures, database tables, and the mapping of inputs to outputs in machine learning models. 

\section{Number Systems}

Before we can organize items into sets, we need to establish the standard sets of numbers used in mathematical analysis. In data science, understanding number systems helps you decide how to represent features in memory (e.g., integers vs. floating-point numbers) and dictates what kind of statistical models you can apply.

\begin{itemize}
    \item \textbf{Natural Numbers ($\N$):} $\N = \{1, 2, 3, \dots\}$. These are counting numbers. In a dataset, the number of clicks on a website or the number of words in an email would be represented using natural numbers. (Note: Some contexts include $0$).
    \item \textbf{Integers ($\Z$):} $\Z = \{\dots, -2, -1, 0, 1, 2, \dots\}$. Integers include negative values. We often use integers to represent changes in stock prices (up or down by whole amounts) or differences in counts over time.
    \item \textbf{Rational Numbers ($\Q$):} $\Q = \{ \frac{p}{q} \mid p, q \in \Z, q \neq 0 \}$. These are numbers with terminating or repeating decimal representations. 
    \item \textbf{Real Numbers ($\R$):} The set of all rational and irrational numbers (such as $\sqrt{2}, \pi, e$). $\R$ represents the entire continuous line of numbers. In machine learning, continuous features like temperature, weight, or calculated probabilities (between 0 and 1) are treated as real numbers.
\end{itemize}

\textbf{Data Science Relevance:} Knowing the domain of your data is the first step in exploratory data analysis. A machine learning algorithm handles a column of Integers ($\Z$) representing categorical IDs very differently from a column of Real numbers ($\R$) representing house prices.

\begin{videocard}{IITM BS Lecture: W1\_L3 Real \& Complex Numbers}{https://www.youtube.com/watch?v=hz7cuJj17wU}
Official IIT Madras lecture explaining rational and irrational numbers in depth.
\end{videocard}

\section{Set Theory}

A \index{Sets}\textbf{set} is a well-defined, unordered collection of distinct objects. The objects contained within a set are called its \textit{elements} or \textit{members}. Sets are usually denoted by uppercase letters ($A, B, C$) and elements by lowercase letters ($a, b, x, y$).

\begin{definitionbox}{Sets and Subsets}
Let $A$ and $B$ be sets.
\begin{itemize}
    \item \textbf{Element of:} $x \in A$ means $x$ is an element of set $A$.
    \item \textbf{Subset:} $A \subseteq B$ ($A$ is a subset of $B$) if every element of $A$ is also in $B$:
    $$\forall x \, (x \in A \implies x \in B)$$
    \item \textbf{Proper Subset:} $A \subset B$ if $A \subseteq B$ and $A \neq B$. This means $B$ has at least one element that is not in $A$.
    \item \textbf{Empty Set ($\emptyset$):} A set containing no elements. Interestingly, $\emptyset \subseteq S$ for any set $S$.
    \item \textbf{Power Set ($\mathcal{P}(S)$):} The set of all subsets of $S$. If $|S| = n$ (meaning $S$ has $n$ elements), then the power set has $|\mathcal{P}(S)| = 2^n$ elements.
\end{itemize}
\end{definitionbox}

\textbf{Data Science Relevance:} In Python, the `set` data structure is used to quickly remove duplicates from a list (e.g., finding all unique words in a document). A dataset itself can be viewed as a set of data points, and finding a subset (like filtering rows where `age > 30`) is a fundamental data manipulation operation.

\begin{videocard}{IITM BS Lecture: W1\_L4 Set Theory}{https://www.youtube.com/watch?v=8z04uTycZpE}
Official IIT Madras lecture introducing sets, subsets, and fundamental operations.
\end{videocard}

\subsection{Set Operations}

Just as we can add and multiply numbers, we can combine sets using specific operations. For any universal set $U$ (the set containing everything under consideration) and subsets $A, B \subseteq U$:

\begin{itemize}
    \item \index{Union}\textbf{Union ($A \cup B$):} Elements in $A$ or $B$ (or both). Logically equivalent to the "OR" operation. 
    \item \index{Intersection}\textbf{Intersection ($A \cap B$):} Elements in both $A$ and $B$. Logically equivalent to the "AND" operation.
    \item \index{Difference}\textbf{Difference ($A \setminus B$ or $A - B$):} Elements in $A$ but not in $B$. This is like saying "Give me users who bought a laptop ($A$) but did NOT buy a mouse ($B$)."
    \item \index{Complement}\textbf{Complement ($A^c$ or $A'$):} Elements in the universal set $U$ that are not in $A$.
    \item \index{Cartesian Product}\textbf{Cartesian Product ($A \times B$):} The set of all possible ordered pairs where the first element is from $A$ and the second is from $B$: $\{(a, b) \mid a \in A, b \in B\}$.
\end{itemize}

\begin{theorembox}{Laws of Set Algebra}
For any sets $A, B$, and $C$:
\begin{itemize}
    \item \textbf{De Morgan's Laws:}
    $$(A \cup B)^c = A^c \cap B^c$$
    $$(A \cap B)^c = A^c \cup B^c$$
    \item \textbf{Distributive Laws:}
    $$A \cap (B \cup C) = (A \cap B) \cup (A \cap C)$$
    $$A \cup (B \cap C) = (A \cup B) \cap (A \cup C)$$
\end{itemize}
\end{theorembox}

\begin{lstlisting}[language=Python, caption=Data Science: Set Operations in Python]
# Define sets of user IDs who bought laptops vs mice
laptop_buyers = {101, 102, 103, 105}
mouse_buyers = {102, 105, 108, 109}

# Intersection: Users who bought BOTH
both = laptop_buyers.intersection(mouse_buyers)
print("Bought both:", both) # {102, 105}

# Difference: Users who bought laptop but NOT mouse (Prime targets for mouse ads)
laptop_only = laptop_buyers.difference(mouse_buyers)
print("Recommend mouse to:", laptop_only) # {101, 103}
\end{lstlisting}

\subsubsection*{Derivation: Proving De Morgan's Law}
Let's prove the first of De Morgan's Laws: $(A \cup B)^c = A^c \cap B^c$. To prove two sets are equal, we must show that any element in the left set is also in the right set, and vice versa.
\begin{enumerate}
    \item Let $x \in (A \cup B)^c$.
    \item By definition of complement, $x \notin (A \cup B)$.
    \item This means $x$ is neither in $A$ nor in $B$.
    \item Logically, this translates to $x \notin A$ \textbf{AND} $x \notin B$.
    \item Therefore, $x \in A^c$ and $x \in B^c$.
    \item By definition of intersection, $x \in (A^c \cap B^c)$.
    \item We have proved $(A \cup B)^c \subseteq A^c \cap B^c$. The reverse logic holds symmetrically, proving equality.
\end{enumerate}

\begin{examplebox}{Pure Math: Power Set Calculation}
Find the power set of $S = \{a, b, c\}$.
\begin{itemize}
    \item Since $|S| = 3$, the power set must contain $2^3 = 8$ elements.
    \item The subsets are:
    \begin{align*}
        \text{0-element subsets:} & \quad \emptyset \\
        \text{1-element subsets:} & \quad \{a\}, \{b\}, \{c\} \\
        \text{2-element subsets:} & \quad \{a,b\}, \{b,c\}, \{a,c\} \\
        \text{3-element subsets:} & \quad \{a,b,c\}
    \end{align*}
    \item Thus, the power set is:
    $$\mathcal{P}(S) = \Big\{ \emptyset, \{a\}, \{b\}, \{c\}, \{a,b\}, \{b,c\}, \{a,c\}, \{a,b,c\} \Big\}$$
\end{itemize}
\end{examplebox}

\begin{examplebox}{Data Science: Sets in Recommendation Systems}
Let universal set $U$ be the catalog of movies on a streaming platform. Let $A$ be the set of movies watched by User 1, and $B$ be the set of movies watched by User 2. 
\begin{itemize}
    \item $A \cap B$ represents movies both users have watched. If this intersection is large, a recommendation system might infer the users have similar tastes.
    \item $A \setminus B$ represents movies User 1 has seen but User 2 hasn't. These movies are prime candidates to recommend to User 2!
\end{itemize}
\end{examplebox}

\section{Relations}

A binary relation $R$ from a set $A$ to a set $B$ is a subset of the Cartesian product $A \times B$. If the ordered pair $(a, b) \in R$, we say "$a$ is related to $b$" and write $a R b$. 

Relations help us formalize connections. For instance, in a relational database, tables are linked by relations (like "Customer $A$ placed Order $B$"). 

\begin{definitionbox}{Properties of Relations on a Set $A$}
A relation $R \subseteq A \times A$ (a relation from set A to itself) is:
\begin{itemize}
    \item \index{Reflexive}\textbf{Reflexive:} if $a R a$ for every $a \in A$. (Every element is related to itself).
    \item \index{Symmetric}\textbf{Symmetric:} if $a R b \implies b R a$ for all $a, b \in A$. (If $A$ is related to $B$, $B$ is also related to $A$).
    \item \index{Transitive}\textbf{Transitive:} if $(a R b \land b R c) \implies a R c$ for all $a, b, c \in A$.
    \item \index{Equivalence Relation}\textbf{Equivalence Relation:} if it is simultaneously reflexive, symmetric, and transitive.
\end{itemize}
\end{definitionbox}

\textbf{Data Science Relevance:} Equivalence relations are vital for grouping and clustering data. If you define a relation "has the same label as," this forms an equivalence relation. It partitions the dataset into distinct, non-overlapping groups (equivalence classes) based on labels.

\begin{examplebox}{Equivalence Relation Proof}
Let $A = \Z$. Define a relation $R$ on $\Z$ by: $x R y \iff x - y$ is divisible by $3$. Prove that $R$ is an equivalence relation.
\begin{enumerate}
    \item \textbf{Reflexivity:} For any $x \in \Z$, $x - x = 0 = 3 \times 0$. Since $0$ is divisible by $3$, $x R x$ is true. Thus, $R$ is reflexive.
    \item \textbf{Symmetry:} If $x R y$, then $x - y = 3k$ for some integer $k$. 
    Multiply by $-1$: $y - x = -3k = 3(-k)$. Since $-k$ is also an integer, $y - x$ is divisible by $3$. Thus, $y R x$, and $R$ is symmetric.
    \item \textbf{Transitivity:} If $x R y$ and $y R z$, then $x - y = 3k$ and $y - z = 3m$ for some integers $k, m$.
    Add the equations: $(x - y) + (y - z) = 3k + 3m \implies x - z = 3(k + m)$.
    Since $k+m$ is an integer, $x - z$ is divisible by $3$. Thus, $x R z$, and $R$ is transitive.
\end{enumerate}
Since $R$ satisfies all three properties, it is an equivalence relation.
\end{examplebox}



\begin{videocard}{IITM BS Lecture: W1\_L6 Relations}{https://www.youtube.com/watch?v=_bNfC5yW6bk}
Official IIT Madras lecture covering cartesian products, equivalence, and partitioning sets.
\end{videocard}

\section{Functions}

A \index{Function}\textbf{function} $f$ from set $A$ to set $B$ (denoted $\func{f}{A}{B}$) is a special type of relation that associates each element $x \in A$ with \textbf{exactly one} element $y \in B$, written as $y = f(x)$. 

\textbf{Data Science Relevance:} The entire goal of supervised machine learning is to approximate an unknown mathematical function $f$ that maps input features $X$ to target outputs $y$ ($y = f(X)$). Understanding function properties helps you understand model expressiveness.

\begin{definitionbox}{Function Terms and Types}
Let $\func{f}{A}{B}$.
\begin{itemize}
    \item \index{Domain}\textbf{Domain:} The set of all valid inputs ($A$).
    \item \index{Codomain}\textbf{Codomain:} The target set containing all possible outputs ($B$).
    \item \index{Range}\textbf{Range/Image:} The set of actual outputs produced by the function: $\text{Im}(f) = \{ f(x) \mid x \in A \}$. Note that the range is a subset of the codomain ($\text{Im}(f) \subseteq B$).
    \item \index{Injective}\textbf{One-to-One (Injective):} $f$ maps distinct inputs to distinct outputs. No two inputs share the same output:
    $$\forall x_1, x_2 \in A \, (f(x_1) = f(x_2) \implies x_1 = x_2)$$
    \item \index{Surjective}\textbf{Onto (Surjective):} The range equals the codomain. Every element in $B$ has at least one pre-image in $A$:
    $$\forall y \in B, \, \exists x \in A \text{ such that } f(x) = y$$
    \item \index{Bijection}\textbf{Bijection (One-to-One Correspondence):} A function that is both injective and surjective. Bijections are perfectly invertible.
\end{itemize}
\end{definitionbox}

\begin{examplebox}{Math \& DS: Injectivity and Surjectivity Analysis}
Analyze the function $\func{f}{\R}{\R}$ defined by $f(x) = 2x + 5$.
\begin{itemize}
    \item \textbf{Injectivity:} Let $f(x_1) = f(x_2)$.
    $$2x_1 + 5 = 2x_2 + 5 \implies 2x_1 = 2x_2 \implies x_1 = x_2$$
    Thus, $f$ is injective. In a data context, this means this transformation is perfectly unique; no two distinct data points will be transformed to the same value.
    \item \textbf{Surjectivity:} Let $y \in \R$ (codomain). We want to find an $x \in \R$ (domain) such that $f(x) = y$.
    $$2x + 5 = y \implies 2x = y - 5 \implies x = \frac{y-5}{2}$$
    Since $y \in \R$, $x = \frac{y-5}{2}$ is a valid real number in the domain.
    Checking: $f(x) = 2\left(\frac{y-5}{2}\right) + 5 = (y - 5) + 5 = y$.
    Thus, $f$ is surjective. 
    \item Since $f$ is both injective and surjective, it is a \textbf{bijection}.
\end{itemize}
\end{examplebox}

\begin{videocard}{IITM BS Lecture: W1\_L7 Functions}{https://www.youtube.com/watch?v=M-nlI2fgaWI}
Official IIT Madras lecture specifically covering domain, codomain, range, and properties of mappings.
\end{videocard}

\section*{Supplementary Lecture Videos}
For further reinforcement and specialized topics, refer to the following official IIT Madras lectures:
\begin{itemize}
    \item \href{https://www.youtube.com/watch?v=F9BZ5JsnjYM}{W1\_L0: Introduction to the course basics}
    \item \href{https://www.youtube.com/watch?v=WEC6jPWvoj8}{W1\_L1: Natural numbers \& their operations}
    \item \href{https://www.youtube.com/watch?v=jHBIJ50DhJQ}{W1\_L2: Rational numbers - fractions \& their properties}
    \item \href{https://www.youtube.com/watch?v=Ue3y-OE_2lE}{W1\_L5: Construction of subsets \& set operations}
    \item \href{https://www.youtube.com/watch?v=bW5O7a1VX4w}{W1\_L5A: Sets examples - membership subsets \& comprehension}
    \item \href{https://www.youtube.com/watch?v=NwYHojyw3SU}{W1\_L5B: Examples of set operations \& counting problems}
    \item \href{https://www.youtube.com/watch?v=6CyagSq3fWA}{W1\_L7A: Relations examples - cartesian products tables \& join operations}
    \item \href{https://www.youtube.com/watch?v=6YrUx0bAnuQ}{W1\_L7B: Function examples - domains ranges \& growth of functions}
    \item \href{https://www.youtube.com/watch?v=98NeCye2i3k}{W1\_L8: Prime numbers - euclids proof infinitude \& applications in cryptography}
    \item \href{https://www.youtube.com/watch?v=xUZeskT6WiQ}{W1\_L9: Why is a number irrational - proof of sqrt2}
    \item \href{https://www.youtube.com/watch?v=kM9nVJks1a8}{W1\_L10: Sets versus collections - russells paradox \& mathematical foundations}
    \item \href{https://www.youtube.com/watch?v=yjJ_pualJr4}{W1\_L11: Degrees of infinity - countable \& uncountable sets}
\end{itemize}

\newpage
\section{Practice Exercises}
\textit{Try to solve these exercises independently. Detailed step-by-step solutions are available in the Appendix at the end of the book.}

\begin{enumerate}
    \item \textbf{Pure Math:} Let $A = \{1, 2, 3\}$ and $B = \{3, 4, 5\}$. Find $(A \cup B) \setminus (A \cap B)$. This operation is known as the symmetric difference.
    \item \textbf{Data Science:} A dataset contains a set of users who clicked an ad ($C$) and users who bought a product ($B$). Express "Users who clicked the ad but did not buy the product" using set notation.
    \item \textbf{Pure Math:} Determine if the relation $R$ on $\R$ defined by $x R y \iff x < y$ is an equivalence relation. Explain why or why not.
    \item \textbf{Data Science:} You define a relation $R$ among users in a database: User A is related to User B if they live in the same city. Is this an equivalence relation?
    \item \textbf{Pure Math:} Prove whether the function $\func{f}{\Z}{\Z}$ given by $f(x) = x^2$ is injective, surjective, both, or neither.
    \item \textbf{Data Science (Feature Scaling):} A common preprocessing step in data science is Min-Max scaling. Consider a simplified version $\func{g}{[0, 10]}{[0, 1]}$ defined by $g(x) = \frac{x}{10}$. Prove that $g$ is a bijection.
\end{enumerate}
